% Riem_Squared_Definitive.tex
% Definitive determination of |Riem|^2 for CP^2 (Fubini-Study metric)
% Resolves inconsistency between agents and numerical code
% Date: 2026-03-22

\documentclass[11pt]{article}
\usepackage{amsmath,amssymb,amsthm}
\usepackage[margin=1in]{geometry}

\newcommand{\Riem}{\mathrm{Riem}}
\newcommand{\Ric}{\mathrm{Ric}}
\newcommand{\Weyl}{\mathrm{W}}
\newcommand{\vol}{\mathrm{Vol}}
\newcommand{\tr}{\mathrm{tr}}

\title{Definitive Determination of $|\Riem|^2$ for $\mathbb{CP}^2$ \\
       (Fubini--Study Metric)}
\author{DFD Research -- Cross-Reference Audit}
\date{22 March 2026}

\begin{document}
\maketitle

\section{Summary of Result}

\begin{center}
\fbox{\parbox{0.85\textwidth}{
\textbf{Definitive answer:} For $\mathbb{CP}^2$ with the standard Fubini--Study
metric (holomorphic sectional curvature $H=4$, sectional curvatures in $[1,4]$):
\[
  |\Riem|^2 \;=\; R_{ijkl}\,R^{ijkl} \;=\; \boxed{192}
\]
This is verified by \textbf{three independent methods}: explicit tensor contraction,
the Gauss--Bonnet theorem ($\chi=3$), and the Hirzebruch signature theorem ($\tau=1$).
}}
\end{center}

\section{Convention and Setup}

We use the standard Fubini--Study metric on $\mathbb{CP}^2$ with:
\begin{itemize}
  \item Real dimension $n = 4$, complex dimension $m = 2$
  \item Holomorphic sectional curvature $H = c = 4$
  \item Sectional curvatures: $\min = 1$, $\max = 4$
  \item Metric signature $(+,+,+,+)$ (Riemannian)
\end{itemize}

\section{Explicit Curvature Tensor}

For a K\"ahler manifold of constant holomorphic sectional curvature $c$
(Kobayashi--Nomizu, Besse), the Riemann tensor in real notation is:
\begin{equation}\label{eq:Riem}
  R_{ijkl} = \frac{c}{4}\bigl(
    g_{ik}g_{jl} - g_{il}g_{jk}
    + J_{ik}J_{jl} - J_{il}J_{jk}
    + 2\,J_{ij}J_{kl}
  \bigr)
\end{equation}
where $J$ is the complex structure ($J^2 = -I$, $J$ orthogonal and antisymmetric
in orthonormal frame).

\subsection{Derived Curvature Quantities}

From \eqref{eq:Riem} with $c = 4$, $n = 4$:

\begin{align}
  \Ric_{ij} &= R^k{}_{ikj} = \frac{c}{2}(m+1)\,g_{ij} = 6\,g_{ij}
  \quad\text{(Einstein)} \\[4pt]
  R &= g^{ij}\Ric_{ij} = 6n = 24 \\[4pt]
  |\Ric|^2 &= \Ric_{ij}\,\Ric^{ij} = n \times 6^2 = 144
\end{align}

The traceless Ricci tensor vanishes: $E_{ij} = \Ric_{ij} - \frac{R}{n}g_{ij}
= 6g_{ij} - 6g_{ij} = 0$.

\subsection{Volume}

\begin{equation}
  \vol(\mathbb{CP}^2, g_{FS}) = \frac{\pi^2}{2}
\end{equation}
(from the Hopf fibration $S^5(1) \to \mathbb{CP}^2$).

\section{Method 1: Explicit Tensor Contraction}

Computing $|\Riem|^2 = R_{ijkl}\,R^{ijkl}$ from \eqref{eq:Riem} with $c = 4$
in an orthonormal frame requires contracting the six terms generated by squaring the
three bracketed terms. Let
\begin{align}
  A_{ijkl} &= g_{ik}g_{jl} - g_{il}g_{jk}, \\
  B_{ijkl} &= J_{ik}J_{jl} - J_{il}J_{jk}, \\
  C_{ijkl} &= 2\,J_{ij}J_{kl}.
\end{align}

The individual contractions (in $n=4$ dimensions with $J_{ij}J^{ij} = 4$) are:
\begin{align}
  A_{ijkl}A^{ijkl} &= 2n(n-1) = 24, \\
  B_{ijkl}B^{ijkl} &= 24, \\
  C_{ijkl}C^{ijkl} &= 4\,J_{ij}J^{ij}\,J_{kl}J^{kl} = 4 \times 4 \times 4 = 64, \\
  2\,A_{ijkl}B^{ijkl} &= 2 \times 8 = 16, \\
  2\,A_{ijkl}C^{ijkl} &= 2 \times 16 = 32, \\
  2\,B_{ijkl}C^{ijkl} &= 2 \times 16 = 32.
\end{align}

Total: $|A+B+C|^2 = 24 + 24 + 64 + 16 + 32 + 32 = 192$.

With the prefactor $(c/4)^2 = 1$ at $c=4$:
\begin{equation}
  \boxed{|\Riem|^2 = 192}
\end{equation}

Numerical verification by explicit 4-index loop over all $4^4 = 256$ components confirms
this exactly.

\section{Method 2: Gauss--Bonnet Theorem}

The Chern--Gauss--Bonnet theorem in dimension 4 states:
\begin{equation}\label{eq:GB}
  \chi(M^4) = \frac{1}{32\pi^2}\int_{M}
  \bigl(|\Riem|^2 - 4\,|\Ric|^2 + R^2\bigr)\,d\mathrm{vol}
\end{equation}

\textbf{Important:} The coefficient is $\frac{1}{32\pi^2}$, \textbf{not} $\frac{1}{8\pi^2}$.
This is the source of error in \texttt{numerical\_a4.py}.

For $\mathbb{CP}^2$: $\chi = 3$, so:
\begin{align}
  3 &= \frac{1}{32\pi^2}\bigl(192 - 4\times 144 + 576\bigr)\times\frac{\pi^2}{2} \\
    &= \frac{1}{32\pi^2}\times 192 \times \frac{\pi^2}{2} \\
    &= \frac{192}{64} = 3 \quad\checkmark
\end{align}

\section{Method 3: Hirzebruch Signature Theorem}

The signature of $\mathbb{CP}^2$ is $\tau = 1$. The Hirzebruch formula gives:
\begin{equation}
  \tau = \frac{p_1}{3} = \frac{1}{3}\cdot\frac{1}{8\pi^2}
  \int_M \frac{1}{4}\,\varepsilon^{ijkl}\,R_{ij}{}^{ab}\,R_{klab}\,d\mathrm{vol}
\end{equation}

Direct computation of the Pontryagin density via $\frac{1}{4}\varepsilon^{ijkl}R_{ijab}R_{kl}{}^{ab}$
yields the value 48. Then:
\begin{align}
  p_1 &= \frac{1}{8\pi^2}\times 48 \times \frac{\pi^2}{2} = 3, \\
  \tau &= \frac{p_1}{3} = 1 \quad\checkmark
\end{align}

\section{Method 4: Curvature Decomposition}

For a 4-dimensional Einstein manifold ($E = 0$):
\begin{equation}
  |\Riem|^2 = |\Weyl|^2 + \frac{R^2}{6}
\end{equation}

\textbf{Important:} The correct formula has $R^2/6$, \textbf{not} $R^2/24$.
This is a second error in \texttt{numerical\_a4.py}.

The general formula in $n$ dimensions for Einstein manifolds is:
\[
  |\Riem|^2 = |\Weyl|^2 + \frac{2R^2}{n(n-1)}
\]
At $n=4$: $\frac{2R^2}{4\times 3} = \frac{R^2}{6}$.

Computing $|\Weyl|^2 = 96$ (confirmed by explicit Weyl tensor construction), we get:
\begin{equation}
  |\Riem|^2 = 96 + \frac{576}{6} = 96 + 96 = 192 \quad\checkmark
\end{equation}

\section{Curvature Operator Decomposition}

The curvature operator $\mathcal{R}: \Lambda^2 \to \Lambda^2$ as a $6\times 6$
matrix has eigenvalues $\{0, 0, 2, 2, 2, 6\}$, with:
\begin{itemize}
  \item Self-dual block: $\text{diag}(6, 0, 0)$, giving
        $W^+ = \text{diag}(4, -2, -2)$ (after subtracting $R/12 = 2$)
  \item Anti-self-dual block: $\text{diag}(2, 2, 2)$, giving
        $W^- = 0$
\end{itemize}
This confirms $\mathbb{CP}^2$ is \textbf{self-dual} ($W^- = 0$), as expected.

\section{Scaling to Other Normalizations}

Under metric scaling $g \to \lambda\,g$: curvature scales as $1/\lambda$,
$|\Riem|^2$ scales as $1/\lambda^2$, volume scales as $\lambda^2$.

\begin{center}
\begin{tabular}{c|c|c|c|c|c}
\hline
$H$ & $R$ & $|\Ric|^2$ & $|\Riem|^2$ & $|\Weyl|^2$ & $\vol$ \\
\hline
4 (standard) & 24 & 144 & \textbf{192} & 96 & $\pi^2/2$ \\
2 & 12 & 36 & 48 & 24 & $2\pi^2$ \\
1 (unit radius) & 6 & 9 & 12 & 6 & $8\pi^2$ \\
\hline
\end{tabular}
\end{center}

All normalizations satisfy $\chi = 3$ and $\tau = 1$ via the Gauss--Bonnet
and Hirzebruch formulas.

\section{Error Diagnosis}

\subsection{Agent 2 and Agent 4: CORRECT}

Both agents used $|\Riem|^2 = 192$ with the standard $H=4$ Fubini--Study metric.
Agent~4 verified via Gauss--Bonnet. These values are correct.

\subsection{\texttt{numerical\_a4.py}: TWO ERRORS}

The code contains two mathematical errors that partially cancel:

\begin{enumerate}
  \item \textbf{Wrong Gauss--Bonnet coefficient} (line 112): Used the formula
  \[
    \chi = \frac{1}{8\pi^2}\int\bigl(|\Weyl|^2 - 2|E|^2 + R^2/24\bigr)\,d\mathrm{vol}
  \]
  The correct formula is either:
  \[
    \chi = \frac{1}{32\pi^2}\int\bigl(|\Riem|^2 - 4|\Ric|^2 + R^2\bigr)\,d\mathrm{vol}
  \]
  or equivalently (for Einstein manifolds):
  \[
    \chi = \frac{1}{32\pi^2}\int\bigl(|\Weyl|^2 + R^2/6\bigr)\,d\mathrm{vol}
  \]

  \item \textbf{Wrong decomposition formula} (line 122): Used
  $|\Riem|^2 = |\Weyl|^2 + R^2/24$, but the correct 4D Einstein formula is
  $|\Riem|^2 = |\Weyl|^2 + R^2/6$.

  The $R^2/24$ factor appears in the \emph{Gauss--Bonnet integrand} expressed
  in terms of Weyl, \textbf{not} in the algebraic decomposition of $|\Riem|^2$.
  These are different formulas that were conflated.
\end{enumerate}

From the wrong Gauss--Bonnet formula, the code derived $|\Weyl|^2 = 72$ (correct value: 96),
then from the wrong decomposition got $|\Riem|^2 = 72 + 24 = 96$ (correct value: $96 + 96 = 192$).

\subsection{The Value 24}

The value $|\Riem|^2 = 24$ does not correspond to any standard normalization of
$\mathbb{CP}^2$. It may arise from confusing $R = 24$ (scalar curvature) with
$|\Riem|^2$, or from $|\Weyl|^2 = 24$ at the $H=2$ normalization.

\section{Impact on $a_4$ Coefficient}

The Gilkey--Seeley--DeWitt $a_4$ coefficient for the scalar Laplacian on a compact
Riemannian manifold without boundary is:
\begin{equation}
  a_4 = \frac{1}{360}\int_M
  \bigl(5R^2 - 2|\Ric|^2 + 2|\Riem|^2\bigr)\,d\mathrm{vol}
\end{equation}

For $\mathbb{CP}^2$ ($H=4$):
\begin{align}
  a_4^{\text{correct}} &= \frac{1}{360}\bigl(5\times 576 - 2\times 144
                          + 2\times 192\bigr)\times\frac{\pi^2}{2} \\
                       &= \frac{1}{360}\times 2976 \times \frac{\pi^2}{2}
                        = \frac{2976}{720}\,\pi^2
                        = \frac{62}{15}\,\pi^2
                        \approx 40.7780
\end{align}

The erroneous value from \texttt{numerical\_a4.py}:
\begin{align}
  a_4^{\text{wrong}} &= \frac{1}{360}\bigl(5\times 576 - 2\times 144
                         + 2\times 96\bigr)\times\frac{\pi^2}{2} \\
                      &= \frac{1}{360}\times 2784 \times \frac{\pi^2}{2}
                       = \frac{2784}{720}\,\pi^2
                       = \frac{58}{15}\,\pi^2
                       \approx 38.1476
\end{align}

The ratio:
\begin{equation}
  \frac{a_4^{\text{correct}}}{a_4^{\text{wrong}}}
  = \frac{2976}{2784} = \frac{62}{58} = \frac{31}{29} \approx 1.06897
\end{equation}

This is a $\sim 6.9\%$ error in $a_4(\mathbb{CP}^2)$, which propagates into all
downstream quantities including the fine structure constant derivation.

\section{Corrected Values Summary}

\begin{center}
\begin{tabular}{l|c|c}
\hline
\textbf{Quantity} & \textbf{Wrong (code)} & \textbf{Correct} \\
\hline
$|\Riem|^2$ & 96 & 192 \\
$|\Weyl|^2$ & 72 & 96 \\
$a_4(\mathbb{CP}^2)$ & $\frac{58}{15}\pi^2$ & $\frac{62}{15}\pi^2$ \\
\hline
\end{tabular}
\end{center}

All downstream calculations using $|\Riem|^2 = 96$ must be corrected to use
$|\Riem|^2 = 192$ (at the $H=4$ normalization).

\end{document}
